% Shared exercise chunk: l5-lds-em-point
\begin{exercisechunk}
\item \textbf{A posterior-mean approximation to the M-step.}
\label[exercise]{ex:l5-lds-em-point}
\exerciseprofile{2}{2}{2}{\excalculation\quad\exinterpretation}{%
Identify the uncertainty terms lost when the exact conditional distribution is
replaced by its posterior-mean trajectory.}
Use the setting and notation of \cref{ex:l5-lds-em}. Replace
\(\Pi_{k,\mathrm{obs}}\), the exact conditional distribution of the hidden
trajectory given \(z_{1:T,\mathrm{obs}}\), by a point mass at
\((\mu_0,\ldots,\mu_T)\). Derive the resulting update for \(A\) when its
required inverse exists, and identify the conditional covariance terms that
this update omits. Explain why this point-trajectory update is not exact EM in
general and need not inherit EM's likelihood-monotonicity guarantee.

When \(\Pi_{k,\mathrm{obs}}\) has a nonsingular Gaussian density, its mean
\((\mu_0,\ldots,\mu_T)\) is also its unique mode. The substitution therefore
uses the most likely hidden trajectory.
\end{exercisechunk}
